\Cref{tab:ablation} varies one factor at a time on seed~0. Single-seed differences below about two
points should be read against the seed standard deviation of roughly one point (\cref{tab:main}).

\paragraph*{Illumination-diverse training}
Re-rendering the five training scenes with the \emph{same} poses but a single Sun direction (that of
the database) has by far the largest effect of any factor we varied. The frozen Galatea head then
reaches only $\fixIllRone\%$ R@1, $\fixIllDrop$ points below illumination-diverse training, and SALAD
drops the same way ($\salFixIllRone\%$). The single-illumination model is slightly \emph{better} when
the query light matches its training light (C0: $\fixIllCzero\%$ vs.\ $\galSzeroCzero\%$), but on the
three hardest conditions it falls to $\fixIllHardVal\%$, against $\galSzeroHard\%$. Fine-tuning makes
this worse, not better: with a single training Sun the fine-tuned Galatea reaches $\ftFixIllRone\%$
($-\ftFixIllDrop$ points; C0 $\ftFixIllCzero\%$ vs.\ $\ftGalSzeroCzero\%$, hard $\ftFixIllHard\%$ vs.\
$\ftGalSzeroHard\%$) and SALAD $\ftSalFixIllRone\%$. A more flexible model fits the training light more
closely, so backbone adaptation is only useful together with illumination diversity. For lunar VPR,
an illumination sweep in the training data, as provided by renderers such as LunarSim or DEM-based
pipelines, matters far more than the choice of aggregation layer.

\paragraph*{Sinkhorn iterations and clusters}
Changing the number of Sinkhorn iterations ($L=1$: $\LoneRone\%$, $L=10$: $\LtenRone\%$) moves R@1 by
less than seed noise, in line with prediction (P3) and the fast contraction of
\cref{prop:contraction}. Halving the clusters to $m=32$ ($4352$-D) costs nothing
($\mThirtytwoRone\%$), which is attractive for bandwidth-limited rovers. With $m=16$ ($2304$-D), R@1
drops to $\mSixteenRone\%$.
